Interatomic potentials for binary collision simulations are usually written
as screened Coulomb potentials
%
\begin{equation}
   V(r) = \frac{1}{4 \pi \varepsilon_0} \frac{Z_1 Z_2 e^2}{r} 
          \cdot \Phi(\frac{r}{a_{12}}).
   \label{eq:screenedCoulomb}
\end{equation}
%
$Z_1$ and $Z_2$ denote the atomic numbers of the ion and the target atom, 
respectively, $r$ the interatomic separation, and 
$e^2/(4 \pi \varepsilon_0) = 14.4~\mbox{eV\AA}$. The asymptotes of the
trajectories, i.e., the scattering angles and the time integrals, are 
calculated from the interatomic potential or, equivalently, from the
screening function.

In IMSIL, there are three options for calculating the scattering angle
$\theta$ in the center-of-mass system: First, by bicubic interpolation in 
precalculated tables of $\cot(\theta/2)$ \cite{I8730} with the energy and 
impact parameter as independent variables. Such tables 
may be set up, in principle, for any interatomic potential. The scattering 
table file corresponding to the universal ZBL potential \cite{I8512} is 
called \texttt{ZBLuniv}, which is used by default. Tables for some of the 
pair-specific potentials are also provided with the program.
 
Alternatively, the scattering angles may be calculated numerically from
the screening function $\Phi(R)$. The second method is to use screening 
functions read from a file and to interpolate them between the tabulated 
values. Tabulated screening functions, calculated according to ZBL \cite{I8512}
or using the DMol DFT program \cite{I9757}, are provided for (almost) all 
atom pairs. The ZBL data are referred to as \texttt{ZBLspec}, as they provide
atom-pair specific screening functions, to distinguish them from the universal
ZBL potential described below.

The third method of calculating scattering angles is to use 
sum-of-exponentials approximations to the screening function together with
numerical integration. The screening function optionally can have a 
linear term
%
\begin{equation}
   \Phi(R) = \sum_{i=1}^n a_i \cdot \exp(-b_i R) 
             \left[ + c + d \, \frac{R}{R_\mathrm{max}} \right] .
    \label{eq:Phi}
\end{equation}
%
The screening length $a_{12}$ is chosen as
%
\begin{equation}
   a_{12} = \frac{0.468~\mbox{\AA}}{\left(Z_1^p + Z_2^p\right)^q}
   \label{eq:ascr}
\end{equation}
%
The predefined values of the parameters $a_i$, $b_i$, $p$, and $q$ for
three commonly used screening functions are listed in Tab.~\ref{tab:screen}. 
These screening functions can be used for any combination of atoms due to
the screening lengths which depend on $Z_1$ and $Z_2$.
%
\begin{table}
    \caption{Predefined parameters for the screening functions and screening 
             lengths \label{tab:screen}}
    \begin{center}
    \begin{tabular}{cccc}
        ~\\
        \multicolumn{4}{c}{Screening functions (Eq.~\ref{eq:Phi})} \\
        \hline
        & ZBLuniv \cite{I8512} & KrC \cite{I7702} & Moliere \cite{I4701} \\
        \hline
        $a_1$ & 0.18175 & 0.190945 & 0.35 \\
        $a_2$ & 0.50986 & 0.473674 & 0.55 \\
        $a_3$ & 0.28022 & 0.335381 & 0.10 \\
        $a_4$ & 0.02817 &  &  \\
        $b_1$ & 3.19980 & 0.278544 & 0.3 \\
        $b_2$ & 0.94229 & 0.637174 & 1.2 \\
        $b_3$ & 0.40290 & 1.919249 & 6.0 \\
        $b_4$ & 0.20162 &  &  \\
        $a_{12}$ (default) & ZBL & Firsov & Firsov \\
        \hline
    \end{tabular}
    \hfill
    \begin{tabular}{cccc}
        ~\\
        \multicolumn{4}{c}{Screening lengths (Eq.~\ref{eq:ascr})}\\
        \hline
        & ZBL \cite{I8512} & Firsov \cite{I5801} & Lindhard \cite{I6301} \\
        \hline
        p & 0.23 & 1/2 & 2/3 \\
        q & 1    & 2/3 & 1/2 \\
        \hline
    \end{tabular}
    \end{center}
\end{table}

Atom-pair specific screening functions using Eq.~(\ref{eq:Phi}) can be
constructed using tables listing the parameters $a_i$, $b_i$ for (almost)
all $Z_1$-$Z_2$ combinations. They are fits to pair-specific ZBL screening
functions (ZBLspec) or to DMol screening functions (NLH, NLHlin), where
NLHlin uses the linear term and needs $R_\mathrm{max}$ as an input. From
$R_\mathrm{max}$ the parameters $c$ and $d$ are calculated so that 
$\Phi(0)=1$ and $\Phi(R_\mathrm{max})=0$. Such tables may also be provided
by the user.
%
\begin{center}
\begin{tabular}{lll}
   parameter \quad & IMSIL name \qquad\qquad    & in record \\
   \hline
   $Z_1$           & derived from ion name      & -- \\
   $M_1$           & derived from ion name or 
                     \texttt{MASS}              & \texttt{\&ATOMS} \\
   $Z_2$           & derived from material name & -- \\
   $M_2$           & derived from material name or
                     \texttt{MASS}              & \texttt{\&ATOMS} \\
   scattering 
   table file      & \texttt{SCATFILE}          & \texttt{\&SNPAR} \\
   quadrature 
   method          & \texttt{QUAD}              & \texttt{\&SNPAR} \\
   screening function
   coefficients    & \texttt{SCRCOEFFILE}       & \texttt{\&SNPAR} \\
   screening 
   function table  & \texttt{SCRFILE}           & \texttt{\&SNPAR} \\
   screening
   length model    & \texttt{SCRLEN}            & \texttt{\&SNPAR} \\
\end{tabular}
\end{center}

The formats of the \texttt{SCATFILE}, \texttt{SCRCOEFFILE}, and 
\texttt{SCRFILE} files are described in Section~\ref{s:input}. 

